\ExplSyntaxOn

\int_new:N \l__binp_int
\int_set:Nn \l__binp_int {4}
\tl_new:N \l__bina_tl
\tl_new:N \l__binb_tl
\tl_new:N \l__binhw_tl
\int_new:N \l__bina_int
\int_new:N \l__binb_int
\int_new:N \l__binc_int
\int_new:N \l__bind_int
\fp_new:N \l__bina_fp

\DeclareDocumentCommand \BinaryPrecision {m}
{
  \int_set:Nn \l__binp_int {#1}
}

\cs_new_nopar:Npn \int_to_bin:Nn #1#2
{
  \tl_clear:N #1
  \int_set:Nn \l__bina_int {#2}
  \prg_replicate:nn {\l__binp_int}
  {
    \int_set:Nn \l__binb_int {\int_mod:nn {\l__bina_int} {2}}
    \tl_put_left:Nx #1 {\int_use:N \l__binb_int}
    \int_set:Nn \l__bina_int {\int_div_truncate:nn {\l__bina_int} {2}}
  }
}

\cs_new_nopar:Npn \bin_hamming_weight:NN #1#2
{
  \tl_set_eq:NN \l__binhw_tl #2
  \tl_replace_all:Nnn \l__binhw_tl {0} {}
  \int_set:Nn #1 {\tl_count:N \l__binhw_tl}
}

\cs_new_nopar:Npn \bin_flip_one:NNn #1#2#3
{
  \tl_set_eq:NN #2#1
  \prg_replicate:nn {#3 - 1}
  {
    \tl_replace_once:Nnn #2 {1} {a}
  }
  \tl_replace_once:Nnn #2 {1} {0}
  \tl_replace_all:Nnn #2 {a} {1}
}

\DeclareDocumentCommand \HammingWeight {m m}
{
  \int_to_bin:Nn \l__binhw_tl {#2}
  \bin_hamming_weight:NN #1 \l__binhw_tl
}

\DeclareDocumentCommand \HyperCube {m}
{
  \int_set:Nn \l__binp_int {#1}
  \int_zero_new:c {l__hyper_0_int}
  \int_set:cn {l__hyper_0_int} {-1}
  \int_step_inline:nnnn {1} {1} {#1}
  {
    \int_zero_new:c {l__hyper_##1_int}
    \int_set:cn  {l__hyper_##1_int} {\int_use:c {l__hyper_\int_eval:n {##1 - 1}_int} * (#1 - ##1 + 1) / ##1}
  }
  \fp_set:Nn \l__bina_fp {2^{#1}-1}
  \int_set:Nn \l__binc_int {\fp_to_int:N \l__bina_fp}
  \int_step_inline:nnnn {0} {1} {\l__binc_int}
  {
    \int_to_bin:Nn \l__bina_tl {##1}
    \bin_hamming_weight:NN \l__binb_int \l__bina_tl
    \node[every~ hypercube~ label/.try] (hg- \l__bina_tl) at (\int_use:c {l__hyper_\int_use:N \l__binb_int _int},\int_use:N \l__binb_int) {\texttt{\l__bina_tl}};
    \int_incr:c {l__hyper_\int_use:N \l__binb_int _int}
    \int_incr:c {l__hyper_\int_use:N \l__binb_int _int}
  }
  \int_step_inline:nnnn {0} {1} {\l__binc_int}
    {
      \int_to_bin:Nn \l__bina_tl {##1}
      \bin_hamming_weight:NN \l__binb_int \l__bina_tl
      \int_step_inline:nnnn {1} {1} {\l__binp_int}
      {
        \bin_flip_one:NNn \l__bina_tl \l__binb_tl {####1}
        \tl_if_eq:NNTF \l__bina_tl \l__binb_tl
          { % do nothing if nodes are the same
          }
          {
            \draw[every~ hypercube~ edge/.try] (hg- \l__bina_tl) -- (hg- \l__binb_tl);
          }
      }
    }
  }

\ExplSyntaxOff

% Macro to cross out a single node
\newcommand{\CrossOutNode}[1]{%
  \draw[red, thick]
    ($(#1) + (-0.3, 0.3)$)
    --
    ($(#1) + (0.3, -0.3)$);
  \draw[red, thick]
    ($(#1) + (-0.3, -0.3)$)
    --
    ($(#1) + (0.3, 0.3)$);
}

% Macro to cross out a list of nodes
\newcommand{\CrossOutNodes}[1]{%
  \foreach \n in {#1} {%
    \CrossOutNode{\n}%
  }
}
% Macro to draw a circle around a single node
\newcommand{\CircleNode}[2][RoyalBlue]{%
  \draw[#1, thick]
    ($(#2)$)
    circle (0.5);
}

% Macro to circle a list of nodes
\newcommand{\CircleNodes}[2][RoyalBlue]{%
  \foreach \n in {#2} {%
    \CircleNode[#1]{\n}%
  }
}

% Überwachungstabelle der Aufgabe mit 10 Nachrichtenbits und 4 Kontrollbits
% (Spalten 1 bis 14, Zeilen c1 bis c4); Massstab im tikzpicture setzen
\newcommand{\ueberwachungZehnVier}{%
  \fill[blue!25] (0,0) rectangle (15,1);
  \fill[blue!15] (0,-4) rectangle (1,0);
  \fill[gray!8] (1,-4) rectangle (15,0);
  \foreach \i in {1,...,10} {
      \node at (\i+.5,.5) {\(a_{\i}\)};
    }
  \foreach \i in {1,...,4} {
      \pgfmathtruncatemacro{\col}{10+\i}
      \node[text=orange!80!black] at (\col+.5,.5) {\(c_{\i}\)};
      \node[text=orange!80!black] at (.5,.5-\i) {\(c_{\i}\)};
    }
  % Überwachungsbereiche von c1 bis c4, eine Liste pro Zeile
  \foreach \cols [count=\z] in {{1,2,3,5,6,7,11},{1,2,4,5,8,9,12},{1,3,4,6,8,10,13},{2,3,4,7,9,10,14}} {
      \foreach \col in \cols {
          \node at (\col+.5,.5-\z) {\(\times\)};
        }
    }
  % Trennung Nachrichtenbits / Kontrollbits
  \draw[blue!70!black,very thick] (11,-4) -- (11,1);
  \draw[step=1,black!30,thin] (0,-4) grid (15,1);
  \draw[black!65] (0,-4) rectangle (15,1);
}

% Bitwerte des Codeworts über den Spaltentiteln
% #1: 14 Bits als Komma-Liste, #2: Spalte magenta hervorheben (0 = keine)
\newcommand{\ueberwachungWerte}[2]{%
  \node[font=\small\itshape,anchor=east] at (1,1.5) {Wert};
  \foreach \w [count=\col] in {#1} {
      \ifnum\col=#2
        \node[magenta,font=\small\bfseries\ttfamily] at (\col+.5,1.5) {\w};
      \else
        \node[font=\small\ttfamily] at (\col+.5,1.5) {\w};
      \fi
    }
}

% Anzahl Einsen unter den Kreuzen jeder Zeile rechts der Tabelle,
% ungerade Anzahlen (Kontrollbit meldet) magenta
\newcommand{\ueberwachungEinsen}[1]{%
  \node[font=\small\itshape] at (16.4,.5) {Einsen};
  \foreach \n [count=\z] in {#1} {
      \ifodd\n
        \node[magenta] at (16.4,.5-\z) {\textbf{\n}\ \(\boldsymbol{\times}\)};
      \else
        \node[green!50!black] at (16.4,.5-\z) {\n\ \checkmark};
      \fi
    }
}


% command to make sure colorboxes don't create whitespace between code lines
\fboxsep=.6pt

% default settings for listings
\lstset{style=python}

\title[Datenintegrität]{Datenintegrität}

\subtitle{Fehler automatisch erkennen und korrigieren}

\author[Blum, Cyril] % (optional)
{Cyril Blum}

\institute[KLW] % (optional)
{
    Fachschaft Informatik\\
    Kantonsschule im Lee
}

\date[\today] % (optional)

\logo{\includegraphics[width=2cm]{Logos/FreeFlower.pdf}}

\titlegraphic{\vspace{-.7cm}\includegraphics[width=.85\textwidth]{"Grundlagen_Info/06_Datenbanken/Figures/title-emojis"}}

\frame{\titlepage}

\begin{frame}{Rückblick: letztes Mal}
    \begin{itemize}
        \item Beim Übertragen von Daten können Fehler auftreten
        \item Ziel: erkennen, wann eine Nachricht fehlerhaft ist
        \item Idee: jede Bitfolge steht für eine Nachricht
              \begin{table}[H]
                  \centering
                  \begin{tabular}{cc}
                      \textbf{Nachricht} & \textbf{Codewort} \\ \midrule
                      \texttt{Rot}       & \texttt{000}      \\ \midrule

                      \texttt{Blau}      & \texttt{011}      \\ \midrule
                      \texttt{Grün}      & \texttt{101}      \\
                  \end{tabular}
              \end{table}
              \only<1>{Wie viele Fehler erkennt diese Codierung?}
              \only<2>{\textbf{1 Bitfehler} (z.B. \texttt{000} $\to$ \texttt{00\red{1}}, \textbf{ungültig!})}
              \only<3>{Was ist der \textbf{Abstand} dieser Codierung?}
              \only<4>{Der \textbf{Abstand einer Codierung} ist die minimale Distanz zwischen zwei Codewörtern (in diesem Fall 2).}
    \end{itemize}
\end{frame}

\begin{frame}[fragile]{Wichtige Erkenntnis}
    \begin{center}
        \textcolor{red}{\textbf{Eine Codierung ist $k$-fehlererkennend, wenn die Codierung mindestens den Mindest-Abstand $k+1$ hat.}}
    \end{center}
    \only<2>{Z.B.: Die Codierung \texttt{000}, \texttt{011} und \texttt{101} hat den kleinsten Abstand 2, also ist die Codierung 1-fehlererkennend.}
\end{frame}


\begin{frame}
    % same wording as exercise 1.13 in the script; box without counter, since
    % the Beamer exercise counter does not match the script numbering
    \begin{tcolorboxTemplate}[
            title=\faEdit{} \theexercisename,
            parbox=false,
            colback=RoyalBlue!10,
            colframe=RoyalBlue,
            colbacktitle=RoyalBlue]
        Finden Sie für die drei Nachrichten \textbf{Rot}, \textbf{Blau} und \textbf{Grün} eine Codierung aus Bitfolgen der Länge 5, die \textbf{1-fehlerkorrigierend} ist.

        Das bedeutet, wenn ein Fehler vorkommt, weiss man nicht nur, dass die Nachricht fehlerhaft ist, sondern man kann das ursprüngliche Codewort aus der fehlerbehafteten Bitfolge eindeutig rekonstruieren. Was können Sie über den Abstand einer solchen Codierung sagen?

        \begin{table}[H]
            \centering
            \small % keeps the exercise box above the footline
            \begin{tabular}{cc}
                \textbf{Farbe} & \textbf{Codierung} \\\midrule
                Rot            & \texttt{00000}     \\\midrule
                Grün           &                    \\\midrule
                Blau           &                    \\
            \end{tabular}
        \end{table}
    \end{tcolorboxTemplate}
\end{frame}

\begin{frame}[fragile]{Mögliche Lösung}
    \begin{table}[H]
            \centering
            \begin{tabular}{cc}
                \textbf{Farbe} & \textbf{Codierung} \\\midrule
                Rot            & \texttt{00000}     \\\midrule
                Grün           & \texttt{00111}     \\\midrule
                Blau           & \texttt{11100}     \\
            \end{tabular}
        \end{table}
\end{frame}

\begin{frame}[fragile]{1-fehlerkorrigierende Codierungen}
    \begin{itemize}
        \item Ziel: Nicht nur Fehler erkennen, sondern auch \textbf{automatisch korrigieren}
        \item Eine Codierung ist \textbf{1-fehlerkorrigierend}, wenn sie jeden einzelnen Bitfehler beheben kann
        \item \textbf{Voraussetzung:} Der Abstand der Codierung ist mindestens 3
    \end{itemize}
    \begin{center}
        \textcolor{red}{\textbf{Eine Codierung ist $k$-fehlerkorrigierend, wenn der Mindestabstand mindestens $2k+1$ beträgt.}}
    \end{center}
\end{frame}


\begin{frame}{Auftrag}{Skript \enquote{Datenintegrität}}
    \begin{itemize}
        \item \aufgabebox{1.13, 1.15}
        \item \challengebox{1.14} (erst nach den Aufgaben)
    \end{itemize}
\end{frame}